% GENERATED FILE. Edit canonical lessons, then run npm run generate:books.
% canonical-source-hash: fnv1a64:a5c9ffb22d40e7cf
% canonical-lessons: TE-C176-lekapote, TE-C176-ayina, TE-C176-dvara, TE-C176-gurinci, TE-C176-lekunda

\chapter{Otherwise, Even so, Through, About, Without}
\label{ch:te-otherwise-even-so-through-about-without}
\hlchaptermodality{176}

\begin{chapteropening}
\textbf{By the end of this chapter:} \emph{I can say otherwise, even so, through, about and without.}

Complete the last lesson of chapter 176: I can say otherwise, even so, through, about and without.
\end{chapteropening}

\section[\texorpdfstring{\te{లేకపోతే} (lēkapōtē)}{లేకపోతే (lēkapōtē)}]{\te{లేకపోతే} (lēkapōtē) — otherwise, or else}

\label{lesson:TE-C176-lekapote}



\begin{quote}
Before the new one: say the Telugu for dust, then the Telugu for smoke.
\end{quote}

\subsection*{You'll want to know: \te{లేకపోతే}}
\textbf{\te{లేకపోతే}} — \emph{lēkapōtē} — \textquotedblleft{}otherwise, or else\textquotedblright{}.

\begin{grammarlens}[title={What you've built}]
One more word for this chapter.
\end{grammarlens}

\subsection*{Your turn}
\begin{itemize}
\raggedright
  \item \emph{Say it:} \emph{lēkapōtē}
  \item \emph{Say it:} \emph{lēkapōtē}, once more
  \item \emph{Say it:} say \emph{lēkapōtē}
  \item \emph{Recall:} say the Telugu for dust, then the Telugu for smoke, then say \emph{lēkapōtē} again
\end{itemize}

\begin{tcolorbox}[breakable,colback=teal!4,colframe=teal!35!black,title={Before you move on}]
What does \te{లేకపోతే} mean? (\textquotedblleft{}Otherwise, or else\textquotedblright{}.) Say it once more.
\end{tcolorbox}
\section[\texorpdfstring{\te{అయినా} (ayinā)}{అయినా (ayinā)}]{\te{అయినా} (ayinā) — even so, still}

\label{lesson:TE-C176-ayina}



\begin{quote}
Before the new one: say the Telugu for smoke, then the Telugu for otherwise.
\end{quote}

\subsection*{You'll want to know: \te{అయినా}}
\textbf{\te{అయినా}} — \emph{ayinā} — \textquotedblleft{}even so, still\textquotedblright{}.

\begin{grammarlens}[title={What you've built}]
One more word for this chapter.
\end{grammarlens}

\subsection*{Your turn}
\begin{itemize}
\raggedright
  \item \emph{Say it:} \emph{ayinā}
  \item \emph{Say it:} \emph{ayinā}, once more
  \item \emph{Say it:} say \emph{ayinā}
  \item \emph{Recall:} say the Telugu for smoke, then the Telugu for otherwise, then say \emph{ayinā} again
\end{itemize}

\begin{tcolorbox}[breakable,colback=teal!4,colframe=teal!35!black,title={Before you move on}]
What does \te{అయినా} mean? (\textquotedblleft{}Even so, still\textquotedblright{}.) Say it once more.
\end{tcolorbox}
\section[\texorpdfstring{\te{ద్వారా} (dvārā)}{ద్వారా (dvārā)}]{\te{ద్వారా} (dvārā) — through, by means of}

\label{lesson:TE-C176-dvara}



\begin{quote}
Before the new one: say the Telugu for otherwise, then the Telugu for even so.
\end{quote}

\subsection*{You'll want to know: \te{ద్వారా}}
\textbf{\te{ద్వారా}} — \emph{dvārā} — \textquotedblleft{}through, by means of\textquotedblright{}.

\begin{grammarlens}[title={What you've built}]
One more word for this chapter.
\end{grammarlens}

\subsection*{Your turn}
\begin{itemize}
\raggedright
  \item \emph{Say it:} \emph{dvārā}
  \item \emph{Say it:} \emph{dvārā}, once more
  \item \emph{Say it:} say \emph{dvārā}
  \item \emph{Recall:} say the Telugu for otherwise, then the Telugu for even so, then say \emph{dvārā} again
\end{itemize}

\begin{tcolorbox}[breakable,colback=teal!4,colframe=teal!35!black,title={Before you move on}]
What does \te{ద్వారా} mean? (\textquotedblleft{}Through, by means of\textquotedblright{}.) Say it once more.
\end{tcolorbox}
\section[\texorpdfstring{\te{గురించి} (guriñci)}{గురించి (guriñci)}]{\te{గురించి} (guriñci) — about, concerning}

\label{lesson:TE-C176-gurinci}



\begin{quote}
Before the new one: say the Telugu for even so, then the Telugu for through.
\end{quote}

\subsection*{You'll want to know: \te{గురించి}}
\textbf{\te{గురించి}} — \emph{guriñci} — \textquotedblleft{}about, concerning\textquotedblright{}.

\begin{grammarlens}[title={What you've built}]
One more word for this chapter.
\end{grammarlens}

\subsection*{Your turn}
\begin{itemize}
\raggedright
  \item \emph{Say it:} \emph{guriñci}
  \item \emph{Say it:} \emph{guriñci}, once more
  \item \emph{Say it:} say \emph{guriñci}
  \item \emph{Recall:} say the Telugu for even so, then the Telugu for through, then say \emph{guriñci} again
\end{itemize}

\begin{tcolorbox}[breakable,colback=teal!4,colframe=teal!35!black,title={Before you move on}]
What does \te{గురించి} mean? (\textquotedblleft{}About, concerning\textquotedblright{}.) Say it once more.
\end{tcolorbox}
\section[\texorpdfstring{\te{లేకుండా} (lēkuṇḍā)}{లేకుండా (lēkuṇḍā)}]{\te{లేకుండా} (lēkuṇḍā) — without}

\label{lesson:TE-C176-lekunda}



\begin{quote}
Before the new one: say the Telugu for through, then the Telugu for about.
\end{quote}

\subsection*{You'll want to know: \te{లేకుండా}}
\textbf{\te{లేకుండా}} — \emph{lēkuṇḍā} — \textquotedblleft{}without\textquotedblright{}.

\begin{grammarlens}[title={What you've built}]
One more word for this chapter.
\end{grammarlens}

\subsection*{Your turn}
\begin{itemize}
\raggedright
  \item \emph{Say it:} \emph{lēkuṇḍā}
  \item \emph{Say it:} \emph{lēkuṇḍā}, once more
  \item \emph{Say it:} say \emph{lēkuṇḍā}
  \item \emph{Recall:} say the Telugu for through, then the Telugu for about, then say \emph{lēkuṇḍā} again
\end{itemize}

\begin{tcolorbox}[breakable,colback=teal!4,colframe=teal!35!black,title={Before you move on}]
What does \te{లేకుండా} mean? (\textquotedblleft{}Without\textquotedblright{}.) Say it once more.
\end{tcolorbox}
